\section{Logical Analysis Squared}

\begin{quotex}
Logical Analysis Squared = Logical Analysis of a Logical Analysis, from January, 2005. This was a commentary on Chapter 6 of Open Secret by \textbf{Wei Wu Wei} used as the talking points for a regularly meeting group discussing non-duality. 
\end{quotex}

\textbf{Point 1:} Objects are only known as the result of reactions of the senses of sentient beings.

\textbf{Commentary:} The conventional wisdom is that something ``out there", let's say an apple, interacts with our senses—visual, olfactory, tactile, for example—which then cause the nerves to send a signal to the brain which somehow recreates an image or copy of the apple in consciousness. This conception is problematical on several different levels. First of all, in this schema, we never actually see the ``real" apple, we just have the image in our mind. No one can get out of the system to see the real apple out there and the image in here to compare the copy to the original. Therefore, the existence of the real apple out there is merely a surmise, since all we know is the apple in consciousness.

Of course, it gets a little more complex than that. In this model, the nervous system is reacting to various wavelengths of light. If so, how does the mélange of various shades of reds and greens get recognized as an ``apple". If we then take a bite out of it, the smell, taste, and texture in some unknown way combine with the visual sense so that we still recognize an ``apple" in our minds. That means we must first have the concept ``apple" in our minds and then be able to categorize all the sensory data under the concept ``apple".

\textbf{Point 2:} The subject itself is not sensorially cognizable as an object.

\textbf{Commentary:} The subject does not appear in the world of appearances, so it is not known.

\textbf{Point 3:} Metaphysically unacceptable

\textbf{Commentary:} Metaphysics as the ultimate science claims to know what is true always and everywhere, in any possible world. As such, its tenets are certain, unlike physics, whose theories are always in principle falsifiable.

\textbf{Point 4:} There is no valid evidence for the existence of an external world

\textbf{Commentary:} Since what we call the ``external" world is what appears in consciousness, which must be the consciousness of some sentient being, the external world is the sentient being.

\textbf{Note:} This does not mean that there is no objective world; the objective world is precisely what appears in consciousness.

\textbf{Point 5:} There is no factual evidence for the existence of sentient beings

\textbf{Commentary:} Of course, by this standard, only objects can appear in consciousness, not subjects. Therefore, the ``subject" appears in consciousness as a concept. So who is there to be conscious?

\textbf{Point 6:} The definition of consciousness

\textbf{Commentary:} Consciousness can not be defined or demonstrated or pointed out. Metaphysically, consciousness is the manifested aspect of the unmanifested. This is consistent with what we learned from Guenon's Multiple States of Being. So consciousness is beyond all concepts and there is no thing that the word ``consciousness" can refer to. Consciousness is properly basic and is inexplicable in terms of something simpler. Therefore, it is fruitless to ask such questions as why do possibilities manifest?

\textbf{Point 7:} Consciousness as absence

\textbf{Commentary:} Consciousness is not a ``thing". Rather, it is the clearing in which things (beings) appear. In consciousness, things are unconcealed (from unmanifested to manifested) and become present to us. This unconcealing (non-action, wei wu) is not the same as a ``projecting" by the Mind of the world into consciousness. We take this unconcealing as the ``real" world, failing to realize the infinite possibilities that remain concealed; clearly, if something is appearing in one way it is not appearing in myriad other possible ways. We identify with this limited world and maintain its communal existence by immersion in everydayness. Idle chit-chat and the popular media and culture serve to reinforce this limited view. We focus on novelty and fashions and celebrity gossip. Instead of seeing the Infinitude of reality as it is, we are constantly checking to see what ``They" think, whether the ``they" be our family, group, or whatever social unit we identify with.

\textbf{Point 8:} Intuition

\textbf{Commentary:} This is the answer to point 3 — the truth of metaphysics is not something that can be expressed mathematically, scientifically, conceptually, verbally, or logically (although it is logical). It must be grasped all at once in intuition. The phenomenal (apparent) universe is the noumenon manifesting in the consciousness of a self which itself is not an object in the phenomenal universe and so must be the noumenon itself.



\flrightit{Posted on 2006-10-26 by Cologero }
